(1) 温度 $T$ 下两种不同理想气体分别为 $N_A,N_B$ 个粒子，初始占据同一容器内以隔板分开的体积 $V_A,V_B$，$V=V_A+V_B$。除去隔板后在共同体积 $V$ 中自由混合，假定温度不变。求熵增。
(2) 求两气体混合后的分压、总压及化学势变化。
(3) 设可用理想半透膜进行等温可逆分离与混合，求最大可取出的混合功。
(4) 比较两种粒子不可区分且为同种理想气体时的结论，说明为何不能直接套用两种气体的混合熵。
